% !TeX root = ../tesis.tex

\section{Del circuito \texorpdfstring{\HaloTwo}{HaloTwo} de Mithril a Cardano}
\label{sec:mithril-a-cardano-del-circuito-halotwo-de-mithril-a-cardano}

Era necesario cerrar la brecha entre dos entornos que, aun compartiendo la palabra \HaloTwo, imponen restricciones muy distintas. Como se estudió en las Secciones~\ref{sec:mithril-relacion-np-del-circuito-stm} y~\ref{sec:mithril-halotwo-y-complejidad}, el circuito \MidnightSTMCircuit estaba diseñado para producir un argumento sucinto sobre \BLS. Del lado de \Cardano no alcanzaba con que el argumento fuera criptográficamente correcto, sino que debía existir un camino realista para transportar el certificado, reconstruir el \transcript, ejecutar la verificación en Aiken y mantenerse por debajo de los límites simultáneos de \maxTxExUnits y \maxTxSize. Esta sección se enfoca en cómo volver la relación \STMRelation verificable \Onchain en \Cardano. Nos apoyamos en un \fork de \PlutusHaloTwoVerifierGen, en una serie de optimizaciones sobre el verificador Aiken emitido por esa herramienta y en una reorganización del flujo de transacciones del \zkBridge para lograr \StakeDistTx y \MintTx.

\subsection{Necesidad y alcance de la adaptación de \PlutusHaloTwoVerifierGen}
\label{subsec:mithril-a-cardano-necesidad-y-alcance-de-la-adaptacion-de-plutushalotwoverifiergen}

La herramienta original \PlutusHaloTwoVerifierGen fue publicada por IOG para generar verificadores Aiken de circuitos pequeños construidos directamente sobre \texttt{halo2\_proofs}, con ejemplos como multiplicación simple, \textit{lookups} y circuitos de prueba \cite{PlutusHalo2VerifierGenRepo, IOGHalo2PlutusVerifier}. A partir de la \textit{verification key} y de una representación de las expresiones del circuito, la herramienta producía el código del validador correspondiente.

Ese flujo no podía aplicarse directamente a \MidnightSTMCircuit. El circuito utilizado por \Mithril está construido sobre las bibliotecas de \Midnight y emplea tipos, relaciones y pasos de verificación que la herramienta original no podía extraer ni representar. Para salvar esa incompatibilidad desarrollamos un \fork de \PlutusHaloTwoVerifierGen. La adaptación no consistió únicamente en admitir un circuito de mayor tamaño: fue necesario extender el proceso completo de extracción y generación para traducir el verificador de \Midnight a una representación ejecutable en Aiken.

En consecuencia, el \fork modificó las principales capas del proceso de generación:
\begin{itemize}[noitemsep]
	\item agregamos soporte para el circuito \MidnightSTMCircuit, con sus \textit{gadgets}, asignaciones de \textit{witness}, tipos auxiliares y \textit{runtime} específico.
	\item incorporamos un extractor nuevo para relaciones de bibliotecas de \Midnight, junto con una capa de conversión para mapear expresiones, compromisos, evaluaciones, rotaciones y pasos del protocolo a la representación intermedia que el generador ya entendía.
	\item eliminamos el backend Plinth y reorientamos toda la herramienta a Aiken.
	\item el generador pasó a exportar entidades estables consumibles por el \zkBridge. Fija qué bytes representan la prueba, qué valores se usan como \textit{public inputs}, qué parte del acumulador se preserva entre transacciones y cómo se serializa el \ReducedRedeemer que más tarde se volverá a hashear \Onchain.
\end{itemize}

La necesidad de adaptar esta capa se vuelve concreta al comparar los \PCS y protocolos de apertura involucrados. Tanto la herramienta original como el circuito real de \Mithril usan compromisos \KZG \cite{Kate2010}, pero no la misma estrategia de multi-apertura. Los ejemplos del generador utilizan \GWCNineteen, compatible con \texttt{halo2\_proofs}, mientras que las bibliotecas de \Midnight utilizan \HaloTwoMultiOpen. La primera estrategia organiza la agregación alrededor de puntos de evaluación; la segunda agrupa los polinomios por \emph{conjuntos de puntos} y reconstruye de otra forma la apertura agregada. Ambas construcciones y su diferencia algebraica se desarrollan en las subsecciones~\ref{subsec:kzg-gwc19} y~\ref{subsec:kzg-halo2-multiopen} del Apéndice B. Como el verificador resultante tiene una estructura distinta, no podíamos reutilizar sin cambios el extractor, el código generado ni el análisis de costos de los ejemplos originales.

\subsection{Optimizaciones del validador Aiken}
\label{subsec:mithril-a-cardano-optimizaciones-del-validador-aiken}

Una vez obtenido un verificador correcto para \MidnightSTMCircuit, el siguiente obstáculo es económico: pasar a respetar los límites de \maxTxSize y \maxTxExUnits de \Cardano. Parte del costo proviene de la frontera entre la representación serializada de la prueba y las operaciones de grupo ejecutadas \Onchain. Como se explica en la subsección~\ref{subsec:ec-point-compression} del Apéndice A, los puntos llegan codificados como bytes y deben descomprimirse para reconstruir elementos de $\GOne$ antes de poder sumarlos o multiplicarlos por escalares. El validador Aiken emitido por el generador realizaba estas descompresiones, además de multiplicaciones escalares, sumas y pairings sobre \BLS. Incluso antes de considerar el flujo completo del bridge, el verificador standalone del circuito real de \Mithril quedaba fuera del presupuesto de CPU.

Las primeras mejoras se hicieron sobre el código Aiken generado, no sobre el protocolo del \zkBridge. En el caso del \MSM izquierdo se reemplazó la evaluación genérica basada en \texttt{foldl} y listas de \texttt{MSMElement} por una cadena \textit{inline} de llamadas a \texttt{addG1} y \texttt{scaleG1}, eliminando \textit{overhead} de la máquina CEK de Aiken, y también omitimos multiplicaciones escalares redundantes cuando el coeficiente es la identidad. La misma idea se extendió luego al \MSM derecho.

Sobre ese mismo \MSM derecho eliminamos dos fuentes de costo adicionales. La primera era una conversión redundante del punto \texttt{vanishing\_g}: una vez computado como elemento de $\GOne$, el verificador lo serializaba nuevamente como \texttt{ByteArray} y luego volvía a descomprimirlo antes de utilizarlo. Mantener el punto en su representación de grupo evita ese recorrido sin alterar la semántica de la verificación. La segunda era la descompresión duplicada de los puntos $W_i$, utilizados tanto en el lado izquierdo como en el derecho. Reutilizar los elementos ya reconstruidos evita repetir llamadas costosas a \texttt{bls12\_381\_G1\_uncompress}.

También eliminamos términos algebraicamente muertos del cálculo del acumulador del polinomio \textit{vanishing}. En particular, el patrón \texttt{scaleG1(zero, xn)} seguido de una suma con el primer compromiso del bloque es innecesario, porque para todo $s \in \Fr$ valen $s \cdot \OO = \OO$, y $\OO + P = P$. Junto con pequeñas mejoras adicionales, esto produjo ahorros medibles de CPU en todos los circuitos de referencia y permitió que ejemplos pesados del repositorio original, como \texttt{atms\_with\_lookups}, entraran en presupuesto.

Sin embargo, ese conjunto de optimizaciones no resolvió por sí mismo el caso \MidnightSTMCircuit. El verificador de \Mithril mejoró, pero siguió siendo demasiado costoso para una sola transacción. Por eso la solución final no fue simplemente ``emitir mejor código'', sino ver cómo reorganizar el flujo para que el costo pesado de verificar un certificado de \Mithril se pague una sola vez y luego pueda ser consumido por la cadena de \StakeDistribution y por \MintTx sin reejecutar el mismo chequeo.

Aún si el presupuesto de CPU fuera suficiente, cargar en cada transacción los scripts grandes del verificador \STM, del \TxsUpdaterMintingValidator y del \MintingValidator volvía inviable el flujo completo al superar el \maxTxSize de \Cardano. Para esto, aprovechamos los \ReferenceScripts introducidos en el \CIP-33 \cite{CIP0033} para persistir scripts pesados \Onchain y referenciarlos desde las transacciones posteriores en lugar de reenviarlos como parte del \textit{witness set}. Podría ser sólo un detalle de implementación, pero fue estrictamente necesario para que el \zkBridge final entre efectivamente en el modelo \eUTxO de \Cardano. En el flujo concretamente implementado, esa decisión aparece en tres transacciones:
\begin{itemize}[noitemsep]
    \item \PublishPhaseOneReferenceScriptTx,
    \item \PublishMintingTxsUpdaterSpendReferenceScriptTx,
    \item \PublishBridgeMintingReferenceScriptTx.
\end{itemize}

Esta decisión de partir el flujo y publicar scripts pesados responde al conjunto de builtins disponible hoy. Si prospera \CIP-133, que propone una \textit{builtin} nativa de \MSM sobre \BLS \cite{CIP0133, PlutusMSMBenchRepo}, una parte importante del bloque lineal que hoy domina el costo del verificador podría colapsar en una sola operación, reduciendo drásticamente tanto el costo algebraico y el \textit{overhead} CEK, y permitiría hacer la verificación de un certificado de \Mithril sin una transacción de dos fases.

\subsection{Partición de la verificación en dos transacciones}
\label{subsec:mithril-a-cardano-particion-de-la-verificacion-en-dos-transacciones}

Como seguíamos fuera del presupuesto de \maxTxExUnits, debimos partir la validación de un argumento \STM de un certificado de \Mithril en dos transacciones consecutivas. En \MidnightSTMCircuit, el trabajo costoso se concentra en el acumulador derecho (recordemos que estamos en \HaloTwoMultiOpen), es decir, en la combinación de compromisos por \textit{point sets}, potencias de $x_1$ y $x_4$, reconstrucciones interpoladas y el término final que incorpora $f(x_3)$. Ese acumulador es el bloque lineal dominante de \ExUnits.

Hicimos el siguiente modelo algebraico usado para diseñar el \textit{split}. Para cada conjunto de puntos $S_j$, el \textit{verifier} arma una combinación lineal $Q_j$ de compromisos usando potencias de $x_1$, evalúa la parte interpolada $R_j(x_3)$, y luego combina los distintos conjuntos con potencias de $x_4$ para formar el acumulador derecho. Si denotamos por $I_j$ al conjunto de índices de compromisos de $S_j$, entonces partir el verificador equivale a escribir $I_j = I_j^{(1)} \sqcup I_j^{(2)}$ y descomponer (por linealidad) tanto $Q_j$ como $R_j(x_3)$ y el acumulador final en una suma de una parte de fase 1 y una parte de fase 2. La primera transacción no ``verifica media prueba'' en un sentido informal, sino que computa un prefijo algebraicamente bien definido del mismo acumulador que la segunda transacción terminará de cerrar.

Para este diseño, debemos decidir cómo transportar un estado intermedio compacto. En lugar de persistir la prueba completa o de recomputar todo en la segunda transacción, la fase 1 produce un \textit{datum} enriquecido que preserva exactamente lo necesario para retomar el cálculo: un acumulador parcial comprimido, tres desafíos del \transcript y un escalar adicional. En la implementación, esa estructura aparece como \PhaseOneState, mientras que la segunda transacción recibe un \ReducedRedeemer que contiene sólo las contribuciones restantes necesarias para completar el bloque lineal y ejecutar el chequeo final de pairing.

El punto de corte finalmente no se eligió sólo por cálculo estático, sino por medición directa. Aunque el análisis preliminar sugería un corte algo más tardío, al instrumentar por separado \PhaseOneSetupTx y \PhaseTwoVerifyTx vimos que la fase 1 absorbía más trabajo fijo del esperado y que el mejor equilibrio real aparecía antes en la entrada 17 del primer conjunto de puntos. Como los cuatro conjuntos de puntos de \MidnightSTMCircuit tienen tamaños $44, 6, 3, 2$, la implementación final deja a la fase 1 las contribuciones de los primeros 17 puntos del primer conjunto, y a la fase 2 todo lo restante. Esta partición fue la que produjo el mejor balance medido entre ambas transacciones (en poco más de 7 billones de unidades), y lo hizo viable.

En el flujo integrado esta partición no se ejecuta para todo certificado que aparece en el protocolo. Las dos corridas \PhaseOneSetupTx $\rightarrow$ \PhaseTwoVerifyTx que permanecen en el camino \Onchain corresponden a los dominios \texttt{stake\_distribution\_standard} y \texttt{cardano\_transactions}. El certificado génesis de \Mithril cumple otro rol: inicializa la cadena de \StakeDistribution y se autentica directamente en la \textit{minting policy} del NFT de \StakeDistribution mediante la firma génesis \GenesisSignature y la clave \GenesisVerificationKey. Por eso no hay un \ProofReceipt asociado al dominio \texttt{stake\_distribution\_genesis}.

\subsection{Validez del transcript y \textit{binding} entre fases}
\label{subsec:mithril-a-cardano-validez-del-transcript-y-binding-entre-fases}
Partir la verificación en dos transacciones introduce una obligación adicional: la fase 2 no puede aceptar un \ReducedRedeemer sólo porque completa algebraicamente el acumulador iniciado por la fase 1. Debe comprobar, además, que ese \redeemer reducido corresponde al mismo \transcript, a la misma prueba serializada y al mismo \textit{statement} público que la primera transacción empezó a procesar. Si esa continuidad fallara, aparecería una superficie de ataque de empalme: un adversario podría usar una fase 1 válida para fijar un estado intermedio y luego intentar cerrar la fase 2 con contribuciones provenientes de otra prueba, otro certificado o un conjunto incompatible de compromisos.

El diseño cierra esa posibilidad con tres mecanismos complementarios. Primero, \PhaseOneState persiste el acumulador parcial y los desafíos necesarios para que la continuación del verificador sea determinista. Segundo, \texttt{reduced\_hash} ata ese estado intermedio al \ReducedRedeemer exacto que deberá aparecer en fase 2. Tercero, \texttt{statement\_hash} conecta el recibo final con el mensaje firmado por el certificado de \Mithril que luego consumirán \StakeDistribution y \MintTx. En conjunto, estos tres enlaces convierten una verificación partida en dos transacciones en una única verificación lógica del mismo objeto criptográfico.

En la primera transacción, \PhaseOneSetupTx recibe la prueba serializada y las dos entradas públicas del circuito, reconstruye el prefijo barato del verificador y deja persistido un datum de tipo \PhaseOneState dentro del \UTxO que luego consumirá la fase 2. Ese estado contiene \texttt{rhs\_prefix}, que es un acumulador parcial comprimido en $\GOne$; \texttt{reduced\_hash}, que fija el hash esperado del \ReducedRedeemer; los desafíos \texttt{x1}, \texttt{x3} y \texttt{x4}; y el escalar \texttt{v}, ya precomputado. Además, el datum extendido del \UTxO de fase 2 conserva un \texttt{statement\_hash}, que en la familia de circuitos \Lagrange coincide directamente con la segunda entrada pública del circuito, es decir, con el mensaje firmado por el certificado de \Mithril.

Es importante notar qué no se persiste. El desafío $x_2$ no reaparece en el datum porque su información ya quedó absorbida en los cálculos que produce la fase 1: una vez reconstruidos $f(x_3)$ y $v$, la segunda transacción ya no necesita rederivar $x_2$ para cerrar el chequeo. Esta decisión reduce el tamaño del estado intermedio y, al mismo tiempo, hace explícita una propiedad conceptual relevante: el datum no guarda ``todo el transcript'', sino solamente la porción de estado necesaria para volver determinista la continuación de la verificación.

El segundo mecanismo de \textit{binding} es el hash del \ReducedRedeemer. La fase 1 serializa canónicamente ese \redeemer reducido, computa \BLAKETwoSTwoFiveSix y persiste el resultado como \texttt{reduced\_hash}. Ese mismo valor se usa además como nombre del NFT emitido en fase 1 y también del NFT que la fase 2 mintea cuando la verificación final tuvo éxito. La consecuencia es fuerte: el estado intermedio, el activo que identifica el \UTxO de fase 2 y el recibo final de prueba verificada quedan todos ligados al mismo objeto serializado. Si la fase 2 intentara presentar un \ReducedRedeemer distinto, el rehash fallaría antes de ejecutar la parte cara del verificador.

El \textit{binding} final se da entre el \texttt{statement\_hash} y el certificado de \Mithril. La fase 2, cuando termina con éxito, no sólo consume el \UTxO intermedio sino que crea un \ProofReceipt con un datum \texttt{ProofReceiptDatum \{ statement\_hash \}}. Ese valor es el mismo que más tarde leerán los validadores de \StakeDistribution y de \MintTx para comprobar que el certificado presentado lleva exactamente el \textit{signed\_message} que ya fue autenticado por \MidnightSTMCircuit. Además, ese \texttt{statement\_hash} no queda respaldado meramente por estar escrito en un datum: queda respaldado por el hecho de que el \UTxO fue creado junto con el NFT minteado bajo la \texttt{PolicyId} de la fase 2, y los validadores downstream buscan precisamente \ProofReceipt marcados por esa política. Así, el \ProofReceipt no es sólo un portador de datos, sino una evidencia autenticada por los scripts que ejecutaron la segunda mitad del verificador.

\subsection{Conexión de \texorpdfstring{\UTxOs, \StakeDistTx\ y \MintTx}{UTxOs, StakeDistTx y MintTx}}
\label{subsec:mithril-a-cardano-conexion-de-utxos-stakedisttx-y-minttx}

El \ProofReceipt anteriormente generado conecta el verificador \STM con el resto del bridge sólo para los certificados que pasan por la ruta de dos fases. En el subflujo de \StakeDistribution existe primero un caso de arranque: \StakeDistributionGenesisTxTemplate no consume un \ProofReceipt, sino que recibe el \GenesisCertificate y verifica su firma génesis \GenesisSignature contra \GenesisVerificationKey mediante la primitiva Ed25519 disponible en Aiken. Si esa autenticación y las condiciones del certificado son válidas, la transacción mintea el NFT de distribución y crea el primer \StakeDistributionUTxO.

Recién después aparece la ruta estándar de \StakeDistribution. \StakeDistributionStandardTxTemplate consume el recibo del dominio \texttt{stake\_distribution\_standard} junto con el \StakeDistributionUTxO padre, verifica la continuidad entre el certificado padre y el certificado nuevo, y emite el nuevo \StakeDistributionUTxO canónico que el resto del protocolo considerará vigente. De esta manera, el estado persistente del bridge en \ChainB queda formado por una cadena de \UTxOs que parte de un certificado génesis autenticado por \GenesisVerificationKey y luego sólo puede avanzar si cada certificado estándar nuevo está ligado a un \ProofReceipt de verificación \STM ya aceptado \Onchain.

En el camino $\LockTx \rightarrow \MintTx$, la transacción \BridgeMintTxTemplate consume el \ProofReceipt del dominio \texttt{cardano\_transactions}, consume también el \TxsUpdaterUTxO y referencia el \StakeDistributionUTxO vigente (como \ReferenceInput, ver \CIP-31 \cite{CIP0031}). La separación de roles es aquí muy precisa. El \ProofReceipt de \texttt{cardano\_transactions} prueba que un \TransactionSnapshotCertificate ha pasado la verificación \STM y fija el \texttt{statement\_hash} que debe coincidir con el \textit{signed\_message} del certificado. El \StakeDistributionUTxO aporta el certificado padre confiable contra el cual debe encadenarse ese \TransactionSnapshotCertificate. Y el \TxsUpdaterMintingValidator garantiza que la transacción bloqueante de \ChainA no pueda reutilizarse para mintear dos veces, reemplazando la nueva raíz de \SMT en el $\TxsUpdaterUTxO^*$ actualizado.

Recién sobre esa base, el argumento \Groth de pertenencia de la \LockTx dentro del \TransactionSnapshot certificado adquiere el significado que el \zkBridge necesita: argumenta inclusión en un snapshot autenticado por un certificado de \Mithril ya enlazado a la \CertificateChain mantenida \Onchain.
